\section{Experimental setup}
\label{sec:experiments}

\subsection{Cross-well Vc low-data benchmark}

The cross-well benchmark predicts Vc windows of length \(N=64\). Training and
validation windows are built from F02-1, F03-2, and F06-1; F03-4 is held out for
testing. The dense input \(u\) concatenates sonic, density, gamma ray, acoustic
impedance, and \(V_p\) windows. Sparse observations \(b\) are generated by a
coordinate subsampling matrix \(M\) with measurement ratios
\(\rho\in\{0.05,0.10,0.20\}\). The low-data regime is controlled by the sliding
window step: step values \(\{8,16,32\}\) yield approximately
\(\{104,52,27\}\) training windows.

For each step and measurement ratio, results are averaged over seeds
\(\{7,23,41\}\). The CLP-CSGM decoder is trained on standardized training
targets, the ridge prior maps log windows to encoded latents, and \(\lambda\)
is selected by validation RMSE.

\subsection{Real-well F03 GR-only porosity benchmark}

The real-well experiment uses F03-4 AC, GR, and porosity data, but the dense
input is intentionally restricted to GR only. Windows have length \(N=64\) and
are split contiguously by depth: 823 training windows, 274 validation windows,
and 275 test windows. Sparse porosity measurements are simulated by coordinate
subsampling with \(\rho\in\{0.20,0.30,0.40,0.50,0.60\}\). This experiment should
be interpreted as a stress test of sparse-measurement assimilation under weak
dense context, not as a claim that GR-only is the preferred operational
petrophysical configuration.

\subsection{Implementation details}

All target windows are standardized using the training-set mean and standard
deviation before autoencoder training and latent optimization; predictions are
mapped back to the original target scale before metric computation. The
CLP-CSGM generator is a multilayer-perceptron autoencoder with one hidden layer
of width 128 in both encoder and decoder, a latent dimension \(k=16\), ReLU
activations, and a linear output layer. The autoencoder is trained for 200
epochs with Adam, learning rate \(10^{-3}\), weight decay \(10^{-5}\), and batch
size 64.

After autoencoder training, the encoder maps each standardized training target
window to a latent code. The paper-facing prior \(h(u)\) is a Ridge regressor
with \(\alpha=1.0\), fitted from standardized log windows to standardized latent
codes. For the prior-class sensitivity experiment, the MLP prior uses two hidden
layers of width 128, \(L_2\) penalty \(10^{-4}\), learning rate \(10^{-3}\), a
maximum of 500 iterations, and early stopping.

At validation and test time, latent refinement is optimized with Adam for 400
iterations, learning rate \(0.05\), and three restarts. The first restart is
initialized at \(z_0=h(u)\); the remaining restarts add small Gaussian
perturbations. The regularization parameter is selected on the validation split
from
\[
  \lambda\in\{10^{-4},3\cdot10^{-4},10^{-3},3\cdot10^{-3},
  10^{-2},3\cdot10^{-2},10^{-1}\}.
\]
For the measurement-only ablation, \(\lambda=0\) and the latent search is
initialized from zero rather than from \(h(u)\). For the prior-only ablation, no
latent optimization is performed and the prediction is \(G(h(u))\).

\subsection{Baselines and metrics}

The main baselines are:
\begin{itemize}
  \item ML only: supervised MLP using dense logs only.
  \item MLP \([u,b]\): direct MLP using concatenated logs and sparse
  measurements.
  \item PCA \([u,b]\): PCA target representation with an MLP from \([u,b]\) to
  PCA scores.
  \item AE \([u,b]\): autoencoder target representation with an MLP from
  \([u,b]\) to latent codes.
\end{itemize}
The proposed method is reported as CLP-CSGM Ridge. We focus on RMSE, MAE, and
relative \(\ell_2\) error. The previous LFISTA/SIR-CS line is not included in
the main comparison because it is a different proposed method with a fixed-basis
sparse residual prior rather than a conditional generative prior.
